\section{Related Work}
\label{sec:related}
\textbf{Audio temporal grounding.}
Weakly supervised grounding learns temporal correspondences from clip-level captions~\cite{xu2024wstag}. Pengi~\cite{deshmukh2023pengi}, Qwen2-Audio~\cite{chu2024qwen2audio}, and Audio Flamingo~3~\cite{goel2025af3} provide audio-conditioned generation. TimeAudio adds temporal markers and absolute time-aware encoding~\cite{wang2026timeaudio}. SpotSound interleaves timestamps with audio tokens and uses absent-event queries to suppress hallucinated intervals~\cite{sun2026spotsound}, combining presence supervision with timestamp prediction. RelTwin targets the query-specific choice between occurrences when inverse relations both have acoustic support.

\textbf{Compositional audio--language learning.}
CLAP aligns audio and text representations~\cite{wu2023clap}; T-CLAP adds order-reversed descriptions and temporal-focused contrast~\cite{yuan2024tclap}. CompA evaluates event order and attribute binding, while CompA-CLAP trains with composition-aware negatives and modular contrast~\cite{ghosh2024compa}. CoSTALA combines spatio-temporal contrast, local event--caption alignment, and feature consistency~\cite{ren2026costala}. These methods organize compositional audio--text representations. RelTwin instead compares the query-conditioned likelihoods of complete timestamp answers, supervising selection in a grounding model's output space.

\textbf{Hard negatives for temporal grounding.}
For compositional video grounding, SHINE combines generated hard-negative queries with within-video and cross-query saliency ranking~\cite{cheng2024shine}, comparing segment relevance and query compositions. RelTwin's candidates are complete timestamp answers within one recording. Each denotes a real event sequence; reversing the query exchanges their labels while preserving the audio. Sequence likelihoods train this query-specific choice.
